% Synthesis of the two-readings note.
\section{Synthesis}
\label{r5:sec:main}

The four questions of the introduction have the following answers; each is stated with its class, and the sections hold the proofs.

\paragraph{(1) Lagrangians.} For the classical reading the free worldline of the arena gives the six-momentum and the shell
$E^2=c^2|\bp|^2+M^2c^4-c^2|\mathbf q|^2$, and the interior shifted Lagrangian gives $\ddot Y=-\kappa^2Y$ with canonical momentum obeying
$\dot\pi=\kappa G\pi$ (restated, \cref{r1:prop:free,r1:prop:shift}). The Dirac Lagrangian $\tfrac12\psi^{\!\top}(K^\mu\partial_\mu-mT)\psi$ gives
$\gamma^\mu\partial_\mu\psi=-mJ\psi$, its current and its stress tensors (restated and inherited, \cref{r1:prop:dirac}). The energies are
governed by the Maxwell Lagrangian $-\tfrac12\Re q$, whose Hamiltonian is the undivided energy $E_6$ and whose first-order form makes
duality an off-shell symmetry (inherited, \cref{r1:prop:maxwell,r1:prop:E6}). The rest energy $E_0=\tfrac12|q|$, taken as a Lagrangian, is
Lorentz invariant and conformal but not duality invariant in the sense of Gaillard and Zumino, has an indefinite energy, and is singular on
radiation; its natural place is the first-order term of the ModMax family. Neither $E_0$ alone nor any Hamiltonian $aE_6+bE_0$ with
$a^2-b^2\neq1$ is Lorentz invariant (fixed, \cref{r1:thm:L0,r1:prop:ham}).

\paragraph{(2) The two $X$'s.} $X_{\rm cl}$ is a displacement, with converted times on $V^+$ and lengths on $V^-$; $X_{\rm EM}=(\bB,\bE)$ is
the pair labelling the six components of the field two-form. They share the vector space and the Clifford action and differ in units,
natural Lorentz algebra ($\lint$ on a four-plane against $\lG$ on the two-form), canonical conjugates and observables. No linear
map equivariant for the interior Lorentz algebra relates them in either direction, whatever dictionary is chosen (fixed,
\cref{r1:thm:type}). The working note's rule of interpretation is confirmed and strengthened.

\paragraph{(3) The four-momentum.} With a clock $n$ on the temporal triple, $E_n=\R n\oplus V^-$ has signature $(1,3)$; the projection
$P_n=(E/c,\bp)$ of the six-momentum is a four-vector of $\mathrm{SO}^+(1,3)_n$, with $\bp$ conjugate to the three lengths and $E/c$ to the
time read by $n$, and it is conserved for every dynamics independent of $(ct_n,\mathbf x)$ (fixed, \cref{r2:prop:cov,r2:prop:cons}). Of
the three mass shells of the working note, the ambient one is the earlier statement that splitting the temporal energy lowers the mass,
the first-class reduction is the paper's surplus-time constraint with the clock $n$, and the third is constructed (restated and
constructed, \cref{r2:prop:shellI,r2:prop:shellII,r2:prop:shellIII}). The four-dimensional mass read by a clock satisfies
$m_4(n)^2c^4=\mu^2c^4-|n\times\mathbf E|^2$, never above the six-dimensional value (fixed, \cref{r2:prop:rotate}). The field has no linear
covariant map to a four-vector; its four-momentum is the integral of its stress tensor, a four-vector only for an isolated whole, and
the photon's $P=\hbar k$ is not derived in the arena (fixed, inherited and open, \cref{r2:sec:field}). The anti-isometry no-go holds for
every clock and obstructs the identification of the interior momentum with the six-momentum, not the duality of the field
(fixed, \cref{r2:prop:nogo,r2:prop:leak}).

\paragraph{(4) The pair on the octants.} The octants are the joint eigenlines of the pole operators, the Fock basis of $S\cong\Lambda^*(L)$.
The field-octant map places the electric member of the pair on the three octants of weight one and the magnetic member on the three of
weight two, uniquely up to a scale and a quarter turn of duality once rotations and Maxwell's equations are imposed; $\bE$ is then a
one-form and $\bB$ a two-form (fixed and constructed, \cref{r3:thm:phi,r3:prop:maxunique}). Maxwell's four equations are one first-order
equation on $S$, with the two Gauss laws on the two poles and the evolution laws on the field octants (fixed, inherited in substance,
\cref{r3:thm:maxwell}); the spatial derivatives in it multiply the temporal generators, a property of the representation and not a
derivation of space from time. The mass form pairs each electric octant with its magnetic antipode and gives $-\bE\cdot\bB$ (fixed,
\cref{r3:prop:energies}). Clifford multiplication by the field point mixes the two members of the pair and does not realise it
(\cref{r3:prop:poles}). Boosts and the duality rotation act on the octants only through odd operators (obstructed for even operators,
\cref{r3:thm:noeven,r3:prop:odd}). The construction represents the classical field strength; it quantises nothing.

\paragraph{The bridge.} The element shared by the two readings is $T=e_1e_2e_3=nJ(n)$: for every clock it is the Dirac mass form of the
particle reading, it lifts the frame parity, and on the octants $-T$ is the duality of the field reading (fixed, \cref{r4:thm:roles}); it is,
up to scale, the only element that is the clock times a Dirac complex structure for every clock (fixed, \cref{r4:thm:unique}). The field
reading does not see the clock and the particle reading does (fixed, \cref{r4:prop:clockdep}). The Dirac action, the octant representation
of the field and the spin lift of $\lG$ are three inequivalent Lorentz representations on one module, whose joint commutant is the
scalars; no bilinear coupling of a Dirac field to the field octants is covariant, while the trilinear couplings are exactly the Pauli
term and its dual, at every clock (fixed and inherited, \cref{r4:prop:ineq,r4:thm:couplings}). The answer to the question whether $T=nJ(n)$
incorporates consistently into both readings is therefore yes at the level of the module: the two readings keep their Lorentz
representations, and they meet in $T$ and in the Pauli coupling. A covariant Dirac--Maxwell dynamics in the arena, with the minimal coupling
through the potential and the charge $J(n)$, a dynamics that selects the clock, and the quantisation of either reading remain open.
